\documentclass[12pt,a4paper]{article}

\usepackage[utf8]{inputenc}
\usepackage[T1]{fontenc}
\usepackage[brazil]{babel}
\usepackage{lmodern}
\usepackage{microtype}
\usepackage[a4paper,margin=2.5cm]{geometry}
\usepackage{setspace}
\usepackage[table]{xcolor}
\usepackage{booktabs}
\usepackage{tabularx}
\usepackage{array}
\usepackage{enumitem}
\usepackage{listings}
\usepackage{fancyhdr}
\usepackage{float}
\usepackage{hyperref}

\hypersetup{
    colorlinks=true,
    linkcolor=blue!50!black,
    urlcolor=blue!60!black,
    citecolor=blue!50!black,
    pdftitle={Trabalho Prático Parte 2: Replicação e Eleição de Líder},
    pdfauthor={Gabriel Welter, Giovanni Milanez, Henrique Trein e Vicente Isatto}
}

\pagestyle{fancy}
\fancyhf{}
\lhead{INF01085 --- Sistemas Distribuídos e Tolerantes a Falhas}
\rhead{Trabalho Prático --- Parte 2}
\cfoot{\thepage}
\setlength{\headheight}{15pt}

\setstretch{1.15}
\setlength{\parindent}{1.25cm}
\setlength{\parskip}{0.25em}

\definecolor{codebg}{RGB}{247,247,247}
\definecolor{codecomment}{RGB}{0,110,50}
\definecolor{codekeyword}{RGB}{0,70,150}

\lstdefinestyle{pythonstyle}{
    language=Python,
    backgroundcolor=\color{codebg},
    basicstyle=\ttfamily\small,
    keywordstyle=\color{codekeyword}\bfseries,
    commentstyle=\color{codecomment},
    stringstyle=\color{orange!60!black},
    numbers=left,
    numberstyle=\tiny\color{gray},
    frame=single,
    breaklines=true,
    showstringspaces=false,
    tabsize=4
}

\newcommand{\arquivo}[1]{\texttt{#1}}
\newcommand{\mensagem}[1]{\texttt{#1}}

\begin{document}

\begin{titlepage}
    \centering
    {\large Universidade Federal do Rio Grande do Sul\par}
    {\large Instituto de Informática\par}
    {\large Departamento de Informática Aplicada\par}
    \vspace{2.5cm}

    {\Large INF01085 --- Sistemas Distribuídos e Tolerantes a Falhas\par}
    \vspace{1.5cm}
    {\huge\bfseries Trabalho Prático --- Parte 2\par}
    \vspace{0.4cm}
    {\LARGE Replicação Ativa e Eleição de Líder\par}

    \vfill
    \begin{tabular}{c}
        Gabriel Weingartner Welter \\
        Giovanni Milanez \\
        Henrique Trein \\
        Vicente Isatto
    \end{tabular}

    \vfill
    Professor Weverton Cordeiro

    \vspace{1cm}
    Porto Alegre, 2026/2
\end{titlepage}


\section{Introdução}

Nesta segunda etapa do trabalho prático, o serviço distribuído de soma desenvolvido na Parte 1 foi estendido para incorporar replicação ativa de máquinas de estados e tolerância a falhas com eleição de líder. O sistema organiza-se em torno de dois papéis desempenhados dinamicamente pelos servidores: o \textit{front-end} (FE), que atua como ponto único de contato dos clientes e sequenciador global atribuindo um identificador monotônico (\arquivo{seq\_id}) a cada requisição e os \textit{replica managers} (RMs), que mantêm cópias locais do acumulador e executam cada operação na mesma ordem determinística. As respostas só são consolidadas e devolvidas aos clientes quando uma maioria simples de confirmações concordantes dos RMs ativos é atingida (quórum). Em caso de queda do líder, os RMs detectam a inatividade e executam o algoritmo do anel, promovendo o servidor de maior endereço IP como novo front-end.

A base de código da Parte 1 foi preservada e adaptada. No cliente (\arquivo{cliente.py}), mantiveram-se a thread de leitura concorrente, a fila e o protocolo \textit{stop-and-wait}, acrescentando-se um identificador aleatório de 64 bits (\arquivo{client\_id}) e um mecanismo de redescoberta por broadcast após timeouts sucessivos. No protocolo de rede (\arquivo{protocolo.py}), foram preservadas as estruturas da etapa anterior e definidas novas mensagens binárias para disseminação ordenada, confirmações locais, batimentos cardíacos (\textit{heartbeats}), troca de topologia, sincronização de estado e eleição no anel. No servidor (\arquivo{servidor.py}), o laço principal de recepção UDP e a arquitetura modular baseada em threads foram mantidos, adicionando-se a alternância dinâmica entre as funções de FE e RM.

Para garantir que requisições de clientes concorrentes não colidissem nos identificadores locais e pudessem ser retransmitidas sem ambiguidade após uma troca de líder, cada operação passou a ser identificada unicamente pela tupla $\langle client\_id, id\_req \rangle$.

\section{Ambiente de testes}

Para a validação e experimentação prática do sistema distribuído, utilizou-se o seguinte ambiente computacional:

\begin{itemize}[leftmargin=1.5cm]
    \item \textbf{Sistema Operacional:} Linux Ubuntu 22.04 LTS (\textit{kernel} 6.8.0-149-generic, x86\_64, glibc 2.35);
    \item \textbf{Processador:} AMD Ryzen 5 7600 (6 núcleos físicos, 12 \textit{threads}, clock base de 3.8 GHz);
    \item \textbf{Memória RAM:} 32 GB DDR5;
    \item \textbf{Linguagem e Interpretador:} Python 3.10.12 (CPython padrão de 64 bits);
    \item \textbf{Topologias de Rede Testadas:}
    \begin{enumerate}[label=(\alph*)]
        \item \textit{Cenário Loopback:} Instâncias de servidores e clientes executadas simultaneamente na mesma máquina, com alocação dinâmica e transparente de endereços de loopback (\texttt{127.0.0.1}, \texttt{127.0.0.2}, \texttt{127.0.0.3}, etc.), atendendo à restrição de não receber IP por linha de comando;
        \item \textit{Cenário Rede Local:} Instâncias distribuídas em diferentes computadores comunicando-se via interface de rede local em um laboratório da UFRGS.
    \end{enumerate}
\end{itemize}

\clearpage

\section{Evolução arquitetural e refatoração}

Conforme requerido na especificação da Parte 2, a arquitetura desenvolvida na Parte 1 passou por uma evolução para viabilizar a transição de um servidor único para um cluster com replicação ativa, ordenação total e eleição em anel lógico. A Tabela~\ref{tab:refatoracao} resume o mapeamento das alterações realizadas entre as etapas.

\begin{table}[H]
    \centering
    \small
    \caption{Mapeamento da evolução arquitetural entre a Parte 1 e a Parte 2.}
    \label{tab:refatoracao}
    \begin{tabularx}{\textwidth}{>{\bfseries}l X X}
        \toprule
        Módulo / Componente & Implementação na Parte 1 & Refatoração e Expansão na Parte 2 \\
        \midrule
        cliente.py & Stop-and-wait simples com identificador incremental \arquivo{id\_req}. Descoberta estática via broadcast pontual. & Inclusão de \arquivo{client\_id} aleatório de 64 bits para formar a chave unívoca $\langle client\_id, id\_req \rangle$. Retransmissão com redescoberta dinâmica do novo FE sob failover. \\
        \addlinespace
        protocolo.py & Apenas 4 mensagens básicas (\mensagem{DISCOVERY}, \mensagem{DISCOVERY\_RESP}, \mensagem{REQ}, \mensagem{ACK}). & Inclusão de 13 novas mensagens intra-cluster (\mensagem{ORDER\_MULTICAST}, \mensagem{RM\_ACK}, \mensagem{HEARTBEAT}, \mensagem{ELECTION}, \mensagem{COORDINATOR}, \mensagem{STATE\_SYNC}, etc.). Serialização com \arquivo{struct} de 64 bits. \\
        \addlinespace
        servidor.py & Processo estático único com papel fixo, sem tolerância a falhas. Acumulador local simples. & Desacoplamento dinâmico entre papéis FE e RM. Detecção automática de IP de interface. Loops assíncronos de controle intra-cluster e orquestração do anel. \\
        \addlinespace
        front\_end.py & Inexistente (integrado ao servidor monolítico da Parte 1). & Atua como Sequenciador Global (\arquivo{seq\_id}), gerencia quórum da maioria simples e responde aos clientes. \\
        \addlinespace
        replica\_manager.py & Inexistente na Parte 1. & Mantém a máquina de estados local, processa a soma de forma ativa e responde ao FE via \mensagem{RM\_ACK}. \\
        \addlinespace
        buffer\_reordenacao.py & Inexistente na Parte 1. & Gerencia buffer FIFO de ordenação total, retém lacunas, detecta duplicatas e garante consistência sequencial. \\
        \addlinespace
        eleicao\_anel.py & Inexistente na Parte 1. & Algoritmo do anel com salto de nós inativos e precedência pelo maior endereço IP. \\
        \addlinespace
        detector\_falhas.py & Inexistente na Parte 1. & Rastreamento baseado em heartbeats periódicos e limiar de falhas consecutivas. \\
        \bottomrule
    \end{tabularx}
\end{table}

\section{Arquitetura da solução}

\subsection{Papéis dos processos}

Existem somente clientes e servidores. O cliente mantém seu papel durante toda
a execução. Um servidor inicia realizando descoberta no canal interno e se não
encontrar um líder, torna-se FE. Caso contrário, ingressa como RM e recebe a
topologia e o estado consolidados.

O servidor que assume o papel de Front-End atua exclusivamente como sequenciador global e agregador de quórum, não integrando o conjunto de réplicas votantes. O processamento da soma e a evolução do acumulador são delegados unicamente aos Replica Managers de forma independente. Uma operação só é confirmada e respondida ao cliente após atingir a maioria simples de confirmações concordantes dos RMs ativos. Na inicialização do cluster, enquanto o primeiro RM ainda não concluiu seu ingresso, o quórum de réplicas permanece indefinido e eventuais requisições aguardam a consolidação da primeira réplica ativa.

\subsection{Canais de comunicação}

A comunicação externa liga clientes ao FE e utiliza descoberta por broadcast,
requisições e respostas unicast. O canal interno liga os servidores e transporta
ordens, confirmações, heartbeats, topologia, eleição e sincronização de estado.

Para a comunicação intra-cluster das ordens de replicação, o método \arquivo{multicast\_rms} realiza o envio direto do datagrama UDP a cada nó registrado na topologia ativa. O foco da camada de replicação é a ordenação total entre os processos por sockets UDP. Essa abordagem dirigida mostra-se mais adequada e previsível para o ambiente de demonstração: em redes de laboratório e switches corporativos, alguns recursos frequentemente bloqueiam ou descartam pacotes de multicast IP na ausência de um roteador de multicast configurado. Além disso, o envio direcionado aos nós da topologia viabiliza a execução de testes em múltiplas interfaces de loopback (\texttt{127.0.0.1}, \texttt{127.0.0.2}, etc.) em uma mesma máquina, assegurando funcionamento confiável sem dependências da infraestrutura física da rede local.

\subsection{Mensagens principais}

\begin{table}[H]
    \centering
    \caption{Principais mensagens do protocolo.}
    \label{tab:mensagens}
    \begin{tabularx}{\textwidth}{>{\ttfamily}l X}
        \toprule
        \normalfont\bfseries Mensagem & \normalfont\bfseries Finalidade \\
        \midrule
        MSG\_DISCOVERY / RESP & Descoberta do FE pelo cliente. \\
        MSG\_REQ / ACK & Requisição de soma e resposta consolidada. \\
        MSG\_ORDER\_MULTICAST & Operação ordenada enviada pelo FE aos RMs. \\
        MSG\_RM\_ACK & Resultado calculado por um RM. \\
        MSG\_NACK & Solicitação de uma sequência ausente. \\
        MSG\_HEARTBEAT & Indicação periódica de atividade e última sequência conhecida. \\
        MSG\_TOPOLOGY & Propagação dos membros e do líder atual. \\
        MSG\_RM\_JOIN / STATE\_SYNC & Ingresso e sincronização de um RM. \\
        MSG\_ELECTION / ELECTION\_ACK & Circulação e confirmação da eleição no anel. \\
        MSG\_COORDINATOR & Proclamação do novo FE. \\
        \bottomrule
    \end{tabularx}
\end{table}

Os inteiros são serializados em ordem de rede por meio de
\arquivo{struct.pack} e \arquivo{struct.unpack}. A identidade de cliente, a
identidade da requisição, a sequência e os acumuladores utilizam 64 bits. IDs de
servidores utilizam 32 bits, o mesmo tamanho de um endereço IPv4.

\section{Replicação ativa}

\subsection{Sequenciamento e idempotência}

O FE mantém um contador global \arquivo{next\_seq} protegido por lock. Para cada
par \arquivo{(client\_id, id\_req)}, ele registra a sequência associada e os
dados da operação. Uma retransmissão encontra o registro existente e não recebe
outra sequência. Assim, a repetição de datagramas do cliente não produz uma nova
soma.

Cada RM utiliza o \arquivo{BufferReordenacao}. Se a sequência recebida é a
esperada, a soma é aplicada. Sequências futuras são retidas até que a lacuna
seja preenchida, e sequências passadas são tratadas como duplicatas. O buffer
preserva temporariamente o resultado original de cada operação e rejeita o uso
do mesmo \arquivo{seq\_id} para conteúdos diferentes.

\subsection{Quórum e consolidação}

O quórum é a maioria simples dos RMs ativos:

\[
    Q = \left\lfloor\frac{M}{2}\right\rfloor + 1,
\]

onde $M$ é a quantidade atual de RMs ativos. Cada RM pode votar uma única vez
por sequência. Os votos são agrupados pelo par
\arquivo{(num\_reqs, total\_sum)}, e o FE responde ao cliente apenas quando um
grupo concordante alcança $Q$. Se um RM falha, ele é removido do conjunto ativo
e o quórum das operações pendentes é reavaliado.

\subsection{Perda de datagramas e sincronização}

Enquanto aguarda quórum, o FE retransmite a ordem com o mesmo
\arquivo{seq\_id}. Quando um RM recebe uma sequência futura, envia
\mensagem{MSG\_NACK} informando a primeira sequência ausente. O FE consulta seu
histórico e reenvia as operações necessárias.

Um RM que ingressa ou retorna recebe \mensagem{MSG\_STATE\_SYNC}, contendo a
próxima sequência, a quantidade de requisições e a soma consolidada. Operações
pendentes também são reenviadas. O RM somente volta ao conjunto ativo quando
seu heartbeat indica que alcançou a sequência do FE. O histórico em memória é
limitado a 10.000 registros para evitar crescimento ilimitado.

\subsection{Garantia de consistência sequencial}

O modelo de replicação ativa implementado baseia-se na Replicação de Máquinas de Estados determinística. O sistema provê a Consistência Sequencial, atendendo às duas condições fundamentais do modelo:

\begin{enumerate}[label=\textbf{\alph*.}]
    \item \textbf{Ordem do Programa:} Para cada cliente individual, as operações são processadas e consolidadas na ordem estrita de emissão. Isso é garantido pelo protocolo de transporte confiável \textit{Stop-and-Wait} presente em \arquivo{cliente.py}, no qual a aplicação aguarda a recepção do \mensagem{MSG\_ACK} com o resultado consolidado da requisição $req_i$ antes de emitir $req_{i+1}$, utilizando a tupla identificadora $\langle client\_id, id\_req \rangle$.
    \item \textbf{Ordem Global Única:} Todas as operações concorrentes de múltiplos clientes são intercaladas e ordenadas em uma sequência global única e monotônica pelo Sequenciador Global (\arquivo{seq\_id} centralizado no Front-End). A disseminação via Multicast Totalmente Ordenado e a barreira do \arquivo{BufferReordenacao} nos Replica Managers asseguram que todas as réplicas processem exatamente a mesma sequência de somas de forma determinística, convergindo para um estado final idêntico do acumulador em todo o cluster.
\end{enumerate}

\section{Detecção de falhas e eleição de líder}

\subsection{Detector de falhas e parametrização temporal}

A monitoração de integridade dos nós opera por meio da troca periódica de mensagens \mensagem{MSG\_HEARTBEAT} via UDP a cada 500~ms, contendo a identidade do líder conhecido e a última sequência processada pelo remetente. O módulo \arquivo{detector\_falhas.py} mantém o registro temporal de contato de cada nó ativo. Em redes locais de baixa latência, adotou-se uma política temporal em dois níveis para balancear a rapidez de recuperação e a tolerância a variações do escalonador do sistema operacional:

\begin{enumerate}[label=\arabic*.]
    \item \textbf{Perda pontual de pacote:} O timeout de espera para retransmissões de requisições de clientes e ordens intra-cluster foi fixado em valores curtos (10~ms a 50~ms), permitindo recuperação rápida diante de eventuais perdas de datagramas UDP sem degradar a vazão.
    \item \textbf{Detecção de falha de nó:} Para evitar eleições desnecessárias causadas por atrasos temporários no agendamento de threads no Python, declara-se a inatividade de um servidor apenas após a ausência de 3 heartbeats consecutivos (timeout acumulado de 1,5~s).
\end{enumerate}

A queda de um RM não inicia eleição: o FE apenas atualiza a topologia e ajusta o quórum necessário das operações em andamento. A ausência de respostas do FE faz com que as réplicas detectem o líder como inativo e disparem imediatamente o Algoritmo do Anel.

\subsection{Algoritmo do anel}

O único mecanismo de eleição adotado na Parte 2 é o \textbf{Algoritmo do Anel} (Chang \& Roberts). A identidade de cada servidor no anel é calculada unicamente a partir de seu endereço IPv4 convertido para um inteiro de 32 bits em ordem de rede, elegendo-se deterministicamente o nó com o maior endereço IP. Os participantes organizam-se logicamente em um anel unidirecional ordenado de forma crescente pelos seus identificadores. O fluxo do algoritmo opera da seguinte forma:
\begin{enumerate}[label=\arabic*.]
    \item Ao constatar o silêncio do Front-End pelo detector de falhas, o nó inicia a eleição emitindo \mensagem{MSG\_ELECTION} ao seu sucessor direto, contendo seu próprio ID como candidato;
    \item Cada repasse no anel requer uma confirmação explícita (\mensagem{MSG\_ELECTION\_ACK}). Caso o sucessor não responda após tentativas consecutivas (timeout de 100~ms), o remetente considera o sucessor inativo, pula-o no percurso e tenta o próximo nó subsequente na topologia, garantindo convergência mesmo diante de quedas sucessivas ou concorrentes;
    \item Ao receber uma mensagem de eleição, o nó compara o candidato recebido com seu próprio identificador: se seu ID for maior, substitui o candidato pelo seu próprio antes de repassar, se for menor, repassa o candidato recebido;
    \item Quando a mensagem retorna com a própria identidade do nó remetente, este torna-se ganhador e autoproclama-se coordenador, mandando \mensagem{MSG\_COORDINATOR} a todos os nós do cluster para assumir as funções de Front-End e sequenciador global.
\end{enumerate}

O novo FE assume com base no seu estado local (próxima sequência, quantidade de requisições e acumulador). Todas as etapas do algoritmo (detecção da queda, disparo da eleição, repasses e proclamação do novo líder) emitem mensagens padronizadas e destacadas no terminal, atendendo às exigências de rastreabilidade/transparência.

\subsection{Redescoberta pelo cliente}

O cliente usa stop-and-wait e mantém a mesma operação até receber o ACK. Após
sucessivas tentativas sem resposta, ele executa novamente a descoberta, atualiza
o endereço do FE e reenvia o mesmo pacote, sem alterar \arquivo{client\_id} ou
\arquivo{id\_req}. Dessa forma, uma troca de líder não transforma a
retransmissão em uma nova operação.

\section{Concorrência e sincronização}

O servidor cria uma thread para cada requisição externa. Há também threads para
o canal de broadcast, o canal interno, heartbeats e monitoramento dos RMs. Locks
separam os estados de eleição, líder, topologia, FE e RM. O uso desses locks
protege o sequenciador, os votos, o conjunto de membros e o acumulador contra
condições de corrida. Operações de rede são realizadas, sempre que possível,
fora das regiões críticas para não bloquear outras requisições.

\section{Problemas encontrados e soluções}

\begin{enumerate}[label=\arabic*.]
    \item \textbf{Idempotência e identificação unívoca de clientes.} Identificar clientes apenas por endereço IP gerava colisões quando múltiplos processos clientes rodavam no mesmo host ou após reinicializações. Introduziu-se um identificador aleatório de 64 bits persistente (\arquivo{client\_id}). A tupla identificadora $\langle client\_id, id\_req \rangle$ é utilizada pelo sequenciador e pelo buffer para detecção e resposta imediata de requisições duplicadas sem duplicar operações no acumulador.

    \item \textbf{Estagnação de réplicas por perda de datagramas UDP.} Caso um datagrama de ordem fosse descartado pela rede local, o RM recebia sequências futuras ($seq\_id > expected\_seq$), retendo pacotes na fila de pendências sem avançar o acumulador. Para evitar travamentos, o RM agora detecta ativamente a lacuna e dispara \mensagem{MSG\_NACK} ao líder, solicitando a retransmissão das fatias ausentes a partir do histórico do Front-End.

    \item \textbf{Salto de sequência e dessincronização no retorno de réplicas.} Ao reiniciar após uma queda temporária, o RM subia com $expected\_seq=1$ enquanto o cluster já operava em sequências adiantadas, retendo permanentemente os pacotes recebidos. Além disso, se o nó retornasse durante uma requisição em andamento no cluster, o envio de $next\_seq$ já incrementado mas com $num\_reqs$ pendente fazia o RM pular a operação em curso. Implementaram-se os métodos \arquivo{obter\_estado\_sincronizacao} e \arquivo{obter\_pacotes\_pendentes} no Front-End: ao receber \mensagem{MSG\_RM\_JOIN}, o líder despacha \mensagem{MSG\_STATE\_SYNC} sincronizando o nó no ponto anterior às pendências e reenvia os pacotes em trânsito para o RM aplicar e votar.

    \item \textbf{Impasse de quórum por divergência de votos e autocura.} Em um cluster com três servidores onde uma réplica caía, o quórum passava a exigir $2/2$ votos concordantes. Caso as duas réplicas restantes possuíssem acumuladores divergentes devido a um retorno defasado no passado, os votos empatavam (1 contra 1), bloqueando o quórum. A solução foi dupla: a reintegração do terceiro RM permite recompor a maioria ($2/3$), e implementou-se um mecanismo de autocura no Front-End que, assim que atinge quórum majoritário, detecta RMs que votaram diferente e despacha \mensagem{MSG\_STATE\_SYNC} corretivo para realinhá-los imediatamente.

    \item \textbf{Reintegração proativa de nós atrasados via Heartbeat.} Réplicas que sofriam partições temporárias de rede voltavam emitindo heartbeats com sequências defasadas em relação ao líder. Em vez de descartar passivamente esses nós, o Front-End no método \arquivo{tratar\_heartbeat} detecta se o remetente relata sequência menor que a esperada e despacha proativamente um \mensagem{MSG\_STATE\_SYNC} para forçar sua atualização e reintegração ao quórum.

    \item \textbf{Continuidade de sequenciamento e desbloqueio de buffer em failover.} Durante a eleição e transição de liderança, dois problemas foram tratados: (a) o novo líder limpa pendências não consolidadas em seu buffer local e assume o sequenciador global a partir da maior sequência conhecida entre os nós ativos; e (b) caso uma operação já tivesse sido aplicada pelas réplicas antes da queda do líder antigo e fosse retransmitida pelo cliente ao novo líder recebendo novo $seq\_id$, as réplicas descartavam o pacote por duplicata. Implementou-se o método \arquivo{resolver\_duplicata\_resequenciada} no \arquivo{BufferReordenacao}, que reconhece a re-sequenciação, avança o $seq\_id$ sem duplicar o valor da soma, libera a fila e devolve \mensagem{MSG\_RM\_ACK} ao novo líder.

    \item \textbf{Ajuste temporal no laço de multicast do Front-End.} O laço de espera por quórum no Front-End aguardava apenas 50 iterações com timeout de 5~ms (total de 250~ms). Quando uma réplica em reinicialização demorava entre 0,5~s e 1~s para ingressar, o FE encerrava o laço sem responder ao cliente, disparando timeouts sucessivos e redescobertas desnecessárias por broadcast. Ampliou-se a janela para 200 iterações (~1~s), concedendo tempo hábil para recuperação de réplicas e consolidação do quórum sob carga.

    \item \textbf{Falha interna de runtime no encerramento forçado (\texttt{SIGINT}).} Ao interromper o servidor com \texttt{Ctrl+C} durante rajadas contínuas de requisições, threads de clientes ainda ativas tentavam imprimir exceções de socket fechado no \texttt{stderr} no mesmo instante em que o interpretador Python finalizava seus descritores de I/O, disparando falha fatal (\textit{\_enter\_buffered\_busy: could not acquire lock for stderr}). O problema foi resolvido tratando silenciosamente os erros de socket fechado nas threads trabalhadoras e invocando \arquivo{os.\_exit(0)} no processo principal, garantindo encerramento limpo pelo sistema operacional.
\end{enumerate}

\section{Testes e resultados experimentais}

A validação do sistema compreendeu testes unitários automatizados, ensaios manuais de tolerância a falhas em ambiente local (\textit{loopback}) e experimentos distribuídos com até cinco computadores na rede local de laboratórios do INF-UFRGS.

\subsection{Testes unitários automatizados}

Por meio do framework \arquivo{unittest}, foram executados os 62 testes da Parte 2 e os 8 testes de regressão da Parte 1, obtendo-se 100\% de aprovação. A Tabela~\ref{tab:testes} sumariza os módulos validados.

\begin{table}[H]
    \centering
    \caption{Resultados da validação do sistema.}
    \label{tab:testes}
    \begin{tabularx}{\textwidth}{X l}
        \toprule
        Cenário / Módulo de Teste & Resultado \\
        \midrule
        Serialização e desserialização binária do protocolo (64 bits) & Aprovado \\
        Ordenação, retenção de lacunas, duplicatas e overflow no buffer & Aprovado \\
        Detector de falhas, heartbeats periódicos e retorno de nós & Aprovado \\
        Eleição em anel (Chang \& Roberts), salto de inativos e convergência & Aprovado \\
        Ingresso dinâmico de nós, sincronização de estado e promoção de FE & Aprovado \\
        Regressão completa das funcionalidades da Parte 1 & Aprovado \\
        Validação manual em loopback (quedas de nós e redescoberta) & Aprovado \\
        Execução distribuída em rede local com 4 e 5 máquinas físicas & Aprovado \\
        Failover com eleição e continuidade do fluxo sob alta carga & Aprovado \\
        Processamento massivo de 1.000.000 de requisições & Aprovado \\
        \bottomrule
    \end{tabularx}
\end{table}

\subsection{Ensaios em ambiente local (Loopback)}

Em uma mesma máquina, instanciaram-se três processos de servidores reutilizando o mesmo par de portas (\texttt{4000} para clientes e \texttt{5000} para o cluster), com alocação automática de endereços de loopback (\texttt{127.0.0.1}, \texttt{127.0.0.2} e \texttt{127.0.0.3}), e um cliente na porta \texttt{4000}. Foram avaliados três cenários fundamentais:

\begin{enumerate}[label=\textbf{\arabic*.}]
    \item \textbf{Queda de Front-End e eleição em anel:} O processo Front-End original (\texttt{127.0.0.1}) foi interrompido com \texttt{Ctrl+C}. As réplicas detectaram a ausência de heartbeats após o limiar configurado e dispararam o Algoritmo do Anel. O nó de maior endereço IP (\texttt{127.0.0.3}) recebeu as confirmações \mensagem{MSG\_ELECTION\_ACK}, autoproclamou-se novo coordenador via \mensagem{MSG\_COORDINATOR} e assumiu o papel de Front-End;
    \item \textbf{Redescoberta e continuidade de requisições pelo cliente:} Após a promoção do novo FE, o cliente emitiu uma requisição de valor 40. Ao esgotar o timeout de espera com o endereço antigo, o cliente executou nova descoberta por broadcast, vinculou-se ao novo Front-End (\texttt{127.0.0.3}) e obteve a confirmação com o acumulador correto (\arquivo{num\_reqs = 1}, \arquivo{total\_sum = 40});
    \item \textbf{Queda de Replica Manager:} A interrupção forçada de um RM ativo não deflagrou eleição indevida. O Front-End registrou o silêncio do nó, removeu-o da topologia ativa e recalculou dinamicamente o quórum para as operações em andamento.
\end{enumerate}

\subsection{Experimentos distribuídos em rede local (Laboratório do INF/UFRGS)}

Para verificar o comportamento do sistema diante de latência real de rede, pacotes concorrentes e múltiplos nós físicos, realizaram-se experimentos no laboratório em duas etapas:

\subsubsection{Cenário com 4 máquinas físicas}

O cluster inicial foi composto por quatro nós em máquinas distintas:
\begin{itemize}
    \item \textbf{Máquina 1:} Inicializada primeiro, detectou ausência de servidores via broadcast e autoproclamou-se Front-End inicial;
    \item \textbf{Máquina 2:} Ingressou como Replica Manager;
    \item \textbf{Máquina 3:} Ingressou como Replica Manager;
    \item \textbf{Máquina 4:} Processo cliente.
\end{itemize}

O cliente conectou-se ao cluster e emitiu uma carga massiva de 1.000.000 de requisições a partir de arquivo texto. No decorrer do envio, o Front-End da Máquina 1 foi derrubado intencionalmente. As réplicas remanescentes detectaram a queda e executaram a eleição no anel. Sendo $\texttt{143.54.50.203} > \texttt{143.54.50.34}$, a Máquina 2 foi eleita e assumiu como novo Front-End. O cliente redescobriu o novo líder e concluiu o lote com sucesso, totalizando:
\[
    num\_reqs = 1.000.000 \quad \text{e} \quad total\_sum = 25.513.505,
\]
com exata equivalência entre o acumulador do cliente e das réplicas ativas.

\noindent\textbf{Diagnóstico e evolução na reintegração de réplicas:} Durante esses ensaios iniciais, observou-se que um nó reiniciado após queda não conseguia reingressar no quórum com o estado atualizado, retendo sequências por aguardar $seq\_id=1$. Notou-se a necessidade de que os nós recuperados deveriam sincronizar o acumulador. Para isso, a equipe desenvolveu e aplicou a sincronização de estado com \mensagem{MSG\_STATE\_SYNC} e descarga de pendências.

\subsubsection{Cenário com 5 máquinas físicas}

Após a implementação da sincronização e recuperação de nós, expandiu-se o cluster para cinco computadores físicos com a adição da Máquina 5. Nessa configuração, o cluster operou com um Front-End e três Replica Managers ativos, estabelecendo quórum da maioria simples de $2/3$. O cliente na Máquina 4 submeteu novamente a carga integral de 1.000.000 de requisições. Foram introduzidas falhas tanto ao FE quanto aos RMs durante o processamento. No final, verificou-se nos terminais que todos os três RMs e o cliente convergiram deterministicamente para o mesmo estado final:
\[
    seq\_id = 1.000.000, \quad num\_reqs = 1.000.000, \quad total\_sum = 25.513.505,
\]
evidenciando a correção da replicação ativa, a estabilidade do quórum sob concorrência e a integridade da máquina de estados replicada.

\section{Uso de ferramentas de inteligência artificial}

Foram utilizados OpenAI Codex e Google Antigravity como ferramentas de apoio técnico durante o desenvolvimento e a validação da Parte 2. Seu uso
concentrou-se na análise comparativa de requisitos com a Parte 1, na escrita e ampliação da suíte de testes unitários, na formatação deste relatório em LaTeX e na depuração de condições de corrida complexas, com destaque para o diagnóstico do bloqueio de buffer sob failover (\arquivo{resolver\_duplicata\_resequenciada}) e o encerramento seguro de threads de I/O via \texttt{SIGINT}.

Todas as sugestões de código e fundamentações teóricas foram revisadas e validadas por meio de testes automatizados e execuções práticas. O uso dos assistentes não substituiu as decisões de projeto do grupo: a arquitetura de replicação ativa, o algoritmo do anel, o modelo de quórum e os mecanismos de recuperação de falhas foram implementados, compreendidos e testados pela equipe.

\section{Divisão de responsabilidades}

A divisão abaixo foi reconstruída a partir do histórico Git do projeto e discussões internas do grupo.

\begin{table}[H]
    \centering
    \caption{Responsabilidades registradas no histórico do projeto.}
    \label{tab:responsabilidades}
    \begin{tabularx}{\textwidth}{>{\bfseries}l X}
        \toprule
        Integrante & Principais contribuições \\
        \midrule
        Gabriel Welter & Organização do estado da Etapa 2, identidade das operações, adaptação inicial do protocolo/cliente, duplicatas e documentação. \\
        Giovanni Milanez & Algoritmo do anel, redescoberta do FE pelo cliente, detector de falhas, heartbeats, failover e correções de concorrência/histórico. \\
        Henrique Trein & Protocolo e buffer iniciais, testes unitários, testes em loopback e rede-local usando múltiplas máquinas, descoberta dinâmica, IP automático, reintegração de RMs, sincronização e correções de failover. \\
        Vicente Isatto & Fluxo FE--RM, disseminação de operações, confirmações, quórum e correções do caminho de replicação. \\
        \bottomrule
    \end{tabularx}
\end{table}

\section{Conclusão}

A implementação desenvolvida atende aos requisitos estabelecidos para a Parte 2: reaproveitamento modular da arquitetura da Parte 1, replicação ativa determinística com quórum por maioria simples, garantia de consistência sequencial via sequenciador centralizado, detecção de falhas baseada em heartbeats periódicos e eleição em anel lógico com salto de nós inativos.

A validação foi feita por testes unitários com cobertura para os casos de borda a experimentos práticos com múltiplos processos concorrentes sob carga real em cenário loopback e rede-local. Os resultados confirmaram a consistência do acumulador mesmo diante de falhas do FE e RMs, e rajadas intensas de requisições, demonstrando a robustez da solução no tratamento de perdas de pacotes, reordenação de mensagens e transições dinâmicas de liderança.

\appendix

\section{Comandos de execução}

Para simplificar a compilação, execução e empacotamento do projeto conforme os requisitos de entrega, disponibilizou-se um \arquivo{Makefile} na raiz do repositório:

\begin{lstlisting}[style=pythonstyle,language=bash,numbers=none]
make compile    # Valida sintaxe e compila bytecodes Python
make test       # Executa os 62 testes unitarios da Etapa 2
make test-all   # Executa os testes de regressao da Etapa 1 e Etapa 2
make zip        # Gera o pacote entrega_etapa2.zip para submissao
\end{lstlisting}

Os servidores escolhem automaticamente o endereço IP e não recebem IP pela linha de comando. Para iniciar várias instâncias na mesma máquina, o mesmo par de portas pode ser reutilizado:

\begin{lstlisting}[style=pythonstyle,language=bash,numbers=none]
python3 Etapa2/servidor.py 4000 5000
python3 Etapa2/servidor.py 4000 5000
python3 Etapa2/servidor.py 4000 5000
\end{lstlisting}

Em outro terminal, inicia-se o cliente:

\begin{lstlisting}[style=pythonstyle,language=bash,numbers=none]
python3 Etapa2/cliente.py 4000
\end{lstlisting}

Para compilar este relatório em uma máquina com distribuição TeX instalada:

\begin{lstlisting}[style=pythonstyle,language=bash,numbers=none]
make report     # ou: cd "Etapa2/Relatorio" && pdflatex main.tex
\end{lstlisting}

\begin{thebibliography}{9}
    \bibitem{enunciado}
    CORDEIRO, Weverton. \textit{Trabalho Prático Parte 2: Replicação e Eleição de Líder}. INF01085 --- Sistemas Distribuídos e Tolerantes a Falhas, UFRGS, 2026/2.

    \bibitem{coulouris}
    COULOURIS, George et al. \textit{Distributed Systems: Concepts and Design}. 5. ed. Addison-Wesley, 2011.
\end{thebibliography}

\end{document}
